\documentclass[11pt]{article}
\usepackage[margin=1.1in]{geometry}
\usepackage[T1]{fontenc}
\usepackage{lmodern}
\usepackage{amsmath,amssymb,amsthm,mathtools}
\usepackage{microtype}
\usepackage{enumitem}
\usepackage[colorlinks=true,linkcolor=black,citecolor=black,urlcolor=black]{hyperref}

\newtheorem{maintheorem}{Theorem}
\renewcommand{\themaintheorem}{\Alph{maintheorem}}
\newtheorem{lemma}{Lemma}
\theoremstyle{definition}
\newtheorem{definition}{Definition}
\theoremstyle{remark}
\newtheorem*{remark}{Remark}

\newcommand{\R}{\mathbb{R}}
\newcommand{\E}{\mathbb{E}}
\newcommand{\cL}{\mathcal{L}}
\newcommand{\eps}{\varepsilon}
\newcommand{\sm}{\setminus}
\newcommand{\ip}[2]{\langle #1,#2\rangle}
\newcommand{\ND}[1]{\textup{(ND#1)}}

\title{Saddle-to-saddle dynamics of deep polynomial ResNets\\[4pt]\large Theorem A: statement and proof}
\author{}
\date{November 2026}

\begin{document}
\maketitle

\noindent
This note states Theorem~A of \texttt{saddle\_to\_saddle.tex} and gives a compressed proof. The result on random initializations (Theorem~B there) is left out. Constants are not tracked: $O(\cdot)$, $C$ and $c$ are uniform on the bounded sets involved. The escape from a saddle (Lemma~\ref{lem:escape}) is proved differently from the long version. There, a contraction estimate in the radial time $\log|w_n|$ gives the escape orbit together with a rate of convergence to it. Here, the escape orbit comes from the centre manifold theorem and, after a polar blow-up, from the stable manifold theorem. The convergence to it comes from a compactness argument, without a rate.

\section{Setting}

\paragraph{Network.} Let $\sigma$ be a polynomial with $\sigma(0)=\dots=\sigma^{(k-1)}(0)=0\ne\sigma^{(k)}(0)$ for some $k\ge3$ (e.g.\ $\sigma(z)=z^k$), and
\[
h_0=x,\qquad h_j=h_{j-1}+\sigma(W_jh_{j-1}+b_j)\quad(1\le j\le L),\qquad f_\theta=h_L,
\]
with $W_j\in\R^{d\times d}$ and $b_j\in\R^d$. A \emph{neuron} $n=(j,i)$ has incoming weights $w_n=(W_j[i,:],b_j[i])\in\R^p$, $p=d+1$. These are all the parameters, since the skip connections play the role of outgoing weights: $\theta=(w_n)_{n\in[N]}\in\Theta=(\R^p)^N$, $N=Ld$. Let $F(\theta)=\cL(f_\theta)$, where $\cL$ is a $C^\infty$ loss on the finite-dimensional space of polynomial maps of degree $\le(\deg\sigma)^L$, e.g.\ $\cL(f)=\frac12\E\|f(x)-y\|^2$. Let $\theta^\alpha$ be the gradient flow $\dot\theta=-\nabla F(\theta)$ with $\theta^\alpha(0)=\alpha\xi$, where $\xi\in\Theta$ is fixed and $\alpha\to0$.

\paragraph{Notation.} For a set $A$ of (active) neurons, $S_A=\{\theta:w_n=0\ \forall n\notin A\}$, $\pi_A$ is the orthogonal projection onto $S_A$, $F_A=F|_{S_A}$, $\nabla_A=\pi_A\nabla$, and $\zeta_A(\theta)=\sum_{n\notin A}|w_n|^k$. We write $\theta^{\sm n}$ for $\theta$ with $w_n$ replaced by $0$, and $\iota_n(w)\in\Theta$ for the parameter with $n$-th component $w$ and all others $0$. The \emph{utility} of $n$ at $\theta$ is the $k$-form
\[
U_n(\theta;w)=-\tfrac1{k!}\,D^k_{w_n}F(\theta^{\sm n})[w,\dots,w],\qquad w\in\R^p,
\]
which depends on $\theta$ only through $\theta^{\sm n}$.

\paragraph{Flatness.} Besides smoothness, the proof uses only
\begin{equation}\tag{H}\label{eq:H}
\partial^\beta_{w_n}F=0\ \text{ on }\{w_n=0\}\qquad\text{for every neuron $n$ and every $1\le|\beta|\le k-1$.}
\end{equation}
Indeed, for $n=(j,i)$ we have $f_\theta=\Phi_j\big(\bar h_j+\sigma(\ip{w_n}{\hat h_{j-1}})e_i\big)$, where $\hat h=(h,1)$, $\bar h_j$ is the output of layer $j$ at $w_n=0$, and $\Phi_j$ is the map computed by the layers after $j$. None of these depends on $w_n$. Since $\sigma(t)=t^k\tau(t)$ for a polynomial $\tau$, every monomial of $f_\theta-f_{\theta^{\sm n}}$ in $w_n$ has degree $\ge k$. Hence $F(\theta)-F(\theta^{\sm n})=\sum_{|\beta|=k}w_n^\beta g_\beta(\theta)$ with $g_\beta$ smooth. Depth, width and biases play no further role: everything below holds for every $F\in C^\infty((\R^p)^N)$ satisfying~\eqref{eq:H}.

\begin{lemma}[Flatness]\label{lem:flat}
Locally uniformly in $\theta$, for every neuron $n$ and every set $A\not\ni n$:
\begin{enumerate}[label=\textup{(\alph*)},nosep]
\item $\nabla_{w_n}F(\theta)=-\nabla U_n(\theta;w_n)+O(|w_n|^k)$; in particular $\nabla_{w_n}F(\theta)=O(|w_n|^{k-1})$;
\item $F(\theta)-F(\pi_A\theta)$ and $\nabla_AF(\theta)-\nabla_AF(\pi_A\theta)$ are $O(\zeta_A(\theta))$, $D_{w_n}\nabla_AF(\theta)=O(|w_n|^{k-1})$, and $\nabla U_n(\theta;w)-\nabla U_n(\pi_A\theta;w)=O(|w|^{k-1}\zeta_A(\theta))$.
\end{enumerate}
\end{lemma}

\begin{proof}
Derivatives of $F$ in directions tangent to $\{w_{n'}=0\}$ still satisfy~\eqref{eq:H} in $w_{n'}$. By Taylor's formula in $w_{n'}$, they change by $O(|w_{n'}|^k)$, with a smooth remainder, when $w_{n'}$ is set to $0$. Setting the $w_{n'}$, $n'\notin A$, to zero one at a time gives (b). For (a), by~\eqref{eq:H} the Taylor polynomial of order $k-1$ of $\nabla_{w_n}F$ in $w_n$ at $w_n=0$ is $-\nabla U_n(\theta;w_n)$.
\end{proof}

\paragraph{Consequences.} By (a), $\nabla_{w_n}F=0$ on $\{w_n=0\}$. So every $S_A$ is invariant under the gradient flow, which on $S_A$ is the gradient flow of $F_A$, and critical points of $F_A$ are critical points of $F$. If $w_n=0$, then $F(\theta+\iota_n(w))=F(\theta)-U_n(\theta;w)+O(|w|^{k+1})$. Hence a critical point of $F_A$ at which some $U_n(\theta;\cdot)$, $n\notin A$, takes a positive value is a saddle of $F$. For the mean squared error, the chain rule gives, with $c_\sigma=\sigma^{(k)}(0)/k!$,
\[
U_{(j,i)}(\theta;w)=c_\sigma\,\E\big[\ip{w}{\hat h_{j-1}}^k\,\big(D\Phi_j(h_j)^\top(y-f_\theta)\big)_i\big].
\]
This is the correlation between the $k$-th power of the would-be pre-activation and the back-propagated residual. At $\theta=0$ it equals $c_\sigma\E[\ip{w}{\hat x}^k(y_i-x_i)]$ for every layer $j$.

\section{The limit process and the result}

A \emph{non-degenerate local minimum} of $F_A$ is a point $\theta^*\in S_A$ with $\nabla F_A(\theta^*)=0$ and $\nabla^2F_A(\theta^*)\succ0$; for $A=\emptyset$, $\theta^*=0$. A \emph{non-degenerate local maximum} of a $k$-form $U$ on $S^{p-1}$ is a critical point $u$ of $U|_{S^{p-1}}$ at which the Riemannian Hessian $(\nabla^2U(u)-kU(u)I)|_{u^\perp}$ is negative definite. Let $\theta^*$ be a non-degenerate local minimum of $F_A$, let $n\notin A$, and let $u^*$ be a non-degenerate local maximum of $U_n(\theta^*;\cdot)|_{S^{p-1}}$ with $U_n(\theta^*;u^*)>0$. The \emph{escape orbit} of $\theta^*$ along $(n,u^*)$ is the maximal gradient-flow trajectory $\gamma$ in $S_{A\cup\{n\}}$ with $\gamma(t)\to\theta^*$ and $\gamma_n(t)/|\gamma_n(t)|\to u^*$ as $t\to-\infty$. It exists and is unique up to time translation (Lemma~\ref{lem:escape}).

\begin{definition}[Saddle-to-saddle process]\label{def:process}
Set $s_0=0$, $A_0=\emptyset$, $\theta^{(0)}=0$ and $\Xi_n(0)=\xi_n$. At stage $m$, we are given $s_m$, $A_m$, a non-degenerate local minimum $\theta^{(m)}$ of $F_{A_m}$, and $\Xi_n(s_m)$ for $n\notin A_m$. Each \emph{dormant} $\Xi_n$, $n\notin A_m$, solves
\begin{equation}\label{eq:dormant}
\tfrac{d}{ds}\Xi_n=\nabla U^{(m)}_n(\Xi_n),\qquad U^{(m)}_n:=U_n(\theta^{(m)};\cdot),
\end{equation}
up to its blow-up time $s_m+T^{(m)}_n\in(s_m,\infty]$. Let $s_{m+1}=s_m+\min_{n\notin A_m}T^{(m)}_n$. Stage $m$ is \emph{non-degenerate} if $s_{m+1}<\infty$ and
\begin{enumerate}[label=\textup{(ND\arabic*)},leftmargin=4em,nosep]
\item the minimum is attained at a single neuron $n_{m+1}$;
\item as $s\uparrow s_{m+1}$, $\Xi_{n_{m+1}}(s)/|\Xi_{n_{m+1}}(s)|$ converges to a non-degenerate local maximum $u_{m+1}$ of $U^{(m)}_{n_{m+1}}|_{S^{p-1}}$;
\item the escape orbit $\gamma_{m+1}$ of $\theta^{(m)}$ along $(n_{m+1},u_{m+1})$ is defined for all $t\ge0$ and converges to a non-degenerate local minimum $\theta^{(m+1)}$ of $F_{A_{m+1}}$, where $A_{m+1}=A_m\cup\{n_{m+1}\}$.
\end{enumerate}
In that case, set $\Xi_n(s_{m+1})=\lim_{s\uparrow s_{m+1}}\Xi_n(s)$ for $n\notin A_{m+1}$ and continue. Otherwise stop, and let $M$ be the last stage reached; $M\le N$, since every jump activates a new neuron. The \emph{limit path} is $\bar\theta(s)=\theta^{(m)}$ for $s\in[s_m,s_{m+1})$, $m\le M$. Let $\bar s_n=s_m$ if $n=n_m$, and $\bar s_n=s_M+T^{(M)}_n$ if $n\notin A_M$.
\end{definition}

\begin{remark}
Write $\Xi_n=R_nu_n$ with $|u_n|=1$, and let $U=U^{(m)}_n$. By Euler's identity $\ip{u}{\nabla U(u)}=kU(u)$, \eqref{eq:dormant} reads $\frac{d}{ds}R_n^{2-k}=-k(k-2)\,U(u_n)$ and $\frac{d}{ds}u_n=R_n^{k-2}\,\nabla_{S}U(u_n)$, where $\nabla_S$ is the spherical gradient. So $n$ activates when its accumulated utility $\int_0^sU^{(m(s'))}_n(u_n)\,ds'$ reaches $|\xi_n|^{2-k}/(k(k-2))$. Meanwhile, its direction performs spherical gradient ascent of the utility. This is the alternation of utility maximization and loss minimization of alternating gradient flows~\cite{KMC25}. In the time $d\tau=R_n^{k-2}ds$, a blow-up forces $\tau\to\infty$ and $U(u_n)\uparrow c>0$. Hence $U^{(m)}_{n_{m+1}}(u_{m+1})>0$ automatically in~\ND2.
\end{remark}

\begin{maintheorem}\label{thm:main}
Let $k\ge3$ and $\xi\in\Theta$. Assume that stages $0,\dots,M-1$ are non-degenerate, and let $S\in(s_M,s_{M+1})$. For $\alpha$ small, $\theta^\alpha$ is defined on $[0,\alpha^{2-k}S]$. As $\alpha\to0$, in the time $s=\alpha^{k-2}t$:
\begin{enumerate}[label=\textup{(\roman*)},nosep]
\item \emph{(plateaus)} for every $\delta>0$, $\theta^\alpha(\alpha^{2-k}s)\to\bar\theta(s)$ uniformly on $[0,S]\sm\bigcup_{1\le m\le M}(s_m-\delta,s_m+\delta)$. In particular the loss converges to the step function $F(\bar\theta(s))$, and $F(\theta^{(0)})>\dots>F(\theta^{(M)})$;
\item \emph{(dormant neurons)} for every neuron $n$ and every $\delta>0$, the rescaled weights $\alpha^{-1}w^\alpha_n(\alpha^{2-k}s)$ converge to $\Xi_n(s)$ uniformly on $[0,\min(S,\bar s_n-\delta)]$;
\item \emph{(jumps)} for $1\le m\le M$, there are times $t^\alpha_m$ with $\alpha^{k-2}t^\alpha_m\to s_m$ such that $\theta^\alpha(t^\alpha_m+\cdot)$ converges to a time translate of $\gamma_m$, locally uniformly on $\R$. Moreover, for $\delta$ smaller than half the gaps between $0,s_1,\dots,s_M,S$, the curve $\{\theta^\alpha(\alpha^{2-k}s):|s-s_m|\le\delta\}$ converges in Hausdorff distance to the closure of $\gamma_m(\R)$.
\end{enumerate}
Each $\theta^{(m)}$ is a critical point of $F$, and a saddle point for $m<M$.
\end{maintheorem}

\section{Proof}

\emph{Idea.} Near a saddle $\theta^*\in S_A$, the active coordinates relax exponentially fast to $\theta^*$. Each inactive weight $w_n=O(\alpha)$ evolves on the slow time scale $\alpha^{2-k}$ by $\dot w_n\approx\nabla U_n(\theta^*;w_n)$ (Lemma~\ref{lem:plateau}). Once a neuron has grown to size $\gg\alpha$ along a non-degenerate maximizer of its utility, it escapes in a vanishing rescaled time while shadowing the escape orbit. The escape is degenerate, since the linearization in $w_n$ vanishes. A centre manifold reduction removes the active coordinates, and in polar coordinates $w_n=\rho u$ the degeneracy can be divided out: the reduced flow has a hyperbolic saddle at $(\rho,u)=(0,u^*)$, whose unstable manifold is the escape orbit (Lemma~\ref{lem:escape}). Theorem~\ref{thm:main} follows by induction on the stages.

\emph{Role of $k\ge3$.} It separates the time scales, since $\alpha^{2-k}\gg1$, and it makes the inactive weights centre directions of each saddle, so that the active coordinates follow them during the escape. It also makes the dormant flows blow up in finite time, at times that depend on $\xi_n$. This breaks the tie between the $L$ neurons of a column, which share the same utility at $\theta=0$. For $k=2$, the dormant flow is linear, a whole column escapes at once, and the outcome of the first jump depends on the ratios of the initial weights.

Throughout, $\theta^*$ is a non-degenerate local minimum of $F_A$, and $\mu>0$ is such that $\nabla^2F_A\succeq\mu$ on $S_A$ near $\theta^*$.

\begin{lemma}[Plateau]\label{lem:plateau}
There are $r_0,C_0>0$ such that the following holds. Let $\Xi_n$, $n\notin A$, solve $\Xi_n'=\nabla U_n(\theta^*;\Xi_n)$ on $[\sigma_0,\sigma_1]$. Let $\tau_\alpha$ be times with $|\pi_A\theta^\alpha(\tau_\alpha)-\theta^*|\le r_0$ and $\alpha^{-1}w^\alpha_n(\tau_\alpha)\to\Xi_n(\sigma_0)$ for $n\notin A$, and let $t(s)=\tau_\alpha+\alpha^{2-k}(s-\sigma_0)$. Then, as $\alpha\to0$, $\alpha^{-1}w_n^\alpha(t(s))\to\Xi_n(s)$ uniformly on $[\sigma_0,\sigma_1]$ for $n\notin A$. Moreover, for small $\alpha$ and $t\in[\tau_\alpha,t(\sigma_1)]$,
\[
|\pi_A\theta^\alpha(t)-\theta^*|\le e^{-\mu(t-\tau_\alpha)}\,|\pi_A\theta^\alpha(\tau_\alpha)-\theta^*|+C_0\sup_{[\tau_\alpha,t]}\zeta_A(\theta^\alpha),\qquad\text{where }\sup\zeta_A(\theta^\alpha)=O(\alpha^k).
\]
\end{lemma}

\begin{proof}
While the rescaled weights $\alpha^{-1}w^\alpha_n$, $n\notin A$, stay bounded, $\zeta_A=O(\alpha^k)$. By Lemma~\ref{lem:flat}(b), $e=\pi_A\theta-\theta^*$ satisfies $\dot e=-\nabla F_A(\theta^*+e)+O(\zeta_A)$, and strong convexity and Grönwall's lemma give the bound. In the time $s$, Lemma~\ref{lem:flat} and the homogeneity of $U_n$ give
\[
\tfrac{d}{ds}\big(\alpha^{-1}w_n\big)=-\alpha^{1-k}\nabla_{w_n}F(\theta)=\nabla U_n\big(\theta^*;\alpha^{-1}w_n\big)+O\big(|\pi_A\theta-\theta^*|+\alpha\big).
\]
By the bound, the error term has integral $O(\alpha^{k-2}+\alpha)$ over $[\sigma_0,\sigma_1]$. Grönwall's lemma gives the convergence, and a continuity argument removes the a priori bound.
\end{proof}

\begin{lemma}[Escape]\label{lem:escape}
Let $n\notin A$, $A'=A\cup\{n\}$ and $U=U_n(\theta^*;\cdot)$, and let $u^*$ be a non-degenerate local maximum of $U|_{S^{p-1}}$ with $U(u^*)>0$. Write $\theta=a+\iota_n(\rho u)+z$ with $a=\pi_A\theta$, $\rho=|w_n|$, $|u|=1$, and $z=\theta-\pi_{A'}\theta$, which collects the other inactive neurons.
\begin{enumerate}[label=\textup{(\alph*)},nosep]
\item The escape orbit $\gamma$ of $\theta^*$ along $(n,u^*)$ exists and is unique up to time translation, and $|\pi_A\gamma-\theta^*|=O(|\gamma_n|^k)$ as $t\to-\infty$. It is the locally uniform limit, as $\eps\to0$, of the gradient-flow trajectories from $\theta^*+\iota_n(\eps u^*)$, shifted so that $|w_n|$ first reaches $\rho_1$ at $t=0$.
\item There are $\rho_1,r,\eps_1,C>0$ with the following property. Suppose that at time $t_0$, $\rho_0:=\rho<\rho_1$, $|u-u^*|\le r$, $\eta:=|z|\le\eps_1\rho_0$ and $|a-\theta^*|\le\rho_0$. Then $\rho$ increases and reaches $\rho_1$ at a time $t_1\le t_0+C\rho_0^{2-k}$, and $|z|\le2\eta$ on $[t_0,t_1]$. Moreover, as $\rho_0\to0$, uniformly in the trajectory, $\theta(t_1)\to\gamma(0)$ and $\sup_{[t_0,t_1]}\operatorname{dist}\big(\theta,\gamma((-\infty,0])\big)\to0$, where $\gamma$ is normalized so that $|\gamma_n|$ first reaches $\rho_1$ at $t=0$.
\end{enumerate}
\end{lemma}

\begin{proof}
Throughout, $r$ is small, and then $\rho_1$ and $\eps_1\le1$ are small.

\emph{Polar coordinates.} By~\eqref{eq:H}, $-\nabla_{w_n}F(a+\iota_n(\rho u)+z)=\rho^{k-1}G(\rho,u,a,z)$, with $G$ smooth for all $\rho\in\R$ and $G(0,u,a,z)=\nabla U_n(a+z;u)$. With $D:=\ip{u}{G}$, the gradient flow reads
\begin{equation}\label{eq:polar}
\dot\rho=\rho^{k-1}D,\qquad\dot u=\rho^{k-2}(G-Du),\qquad\dot a=-\nabla_AF.
\end{equation}
At $(\rho,a,z)=(0,\theta^*,0)$, Euler's identity gives $D=kU(u)$ and $G-Du=\nabla_SU(u)$, the spherical gradient of $U$. Its zero $u^*$ is a sink, with the Riemannian Hessian as linearization.

\emph{Invariance.} Under the assumptions of (b), a continuity argument gives the following on $[t_0,t_1]$. First, $\kappa\le D\le C$ with $\kappa>0$. So $\dot\rho\ge\kappa\rho^{k-1}$: $\rho$ increases, and $t_1-t_0\le C\rho_0^{2-k}$. Next, $\frac{d}{dt}|z|^{2-k}\ge-C$ by Lemma~\ref{lem:flat}(a), so $|z|^{2-k}\ge\eta^{2-k}-C\rho_0^{2-k}\ge(2\eta)^{2-k}$, because $\eta\le\eps_1\rho_0$. Next, $|u-u^*|\le r$: on $\{|u-u^*|=r\}$, $\dot u$ points inwards, since $\nabla_SU$ does and $G-Du-\nabla_SU(u)=O(\rho+|a-\theta^*|+|z|)$ is small. Finally, Lemma~\ref{lem:flat}(b) and the strong convexity of $F_A$ give $\frac{d}{dt}|a-\theta^*|\le-\mu|a-\theta^*|+C(\rho^k+|z|^k)$. Since $\rho$ increases and $|z|\le2\rho$,
\begin{equation}\label{eq:slaving}
|a-\theta^*|\le e^{-\mu(t-t_0)}\rho_0+C\rho^k\qquad\text{on }[t_0,t_1].
\end{equation}

\emph{Centre manifold.} By~\eqref{eq:H} and $k\ge3$, the linearization at $\theta^*$ of the gradient flow of $F_{A'}$ is $-\nabla^2F_A(\theta^*)\oplus0$ on $S_{A'}=S_A\oplus\R^p$. By the centre manifold theorem~\cite[Chapter~2, Theorem~1]{Car81}, this flow has a locally invariant manifold $W^c=\{a=h(w_n)\}$ near $\theta^*$, with $h$ of class $C^2$, $h(0)=\theta^*$ and $Dh(0)=0$. By the invariance of $W^c$, $d:=a-h(w_n)$ satisfies $\dot d=-\big(\nabla^2F_A(\theta^*)+o(1)\big)d$ along trajectories of $S_{A'}$ near $\theta^*$. Hence such a trajectory lies in $W^c$ if it stays near $\theta^*$ for all $t\le t'$.

\emph{Blow-up.} On $W^c$, \eqref{eq:polar} becomes $\dot\rho=\rho^{k-1}\tilde D$, $\dot u=\rho^{k-2}(\tilde G-\tilde Du)$, with $\tilde G(\rho,u):=G(\rho,u,h(\rho u),0)$ and $\tilde D:=\ip{u}{\tilde G}$. In the time $\tau$ defined by $d\tau=\rho^{k-2}dt$, this is the flow of the vector field $X:=(\rho\tilde D,\,\tilde G-\tilde Du)$ on $\R\times S^{p-1}$, which is $C^2$ near $(0,u^*)$. It preserves $\{\rho=0\}$, where it is the spherical gradient flow of $U$. Hence $(0,u^*)$ is a hyperbolic equilibrium of $X$: the eigenvalue $kU(u^*)>0$ has an eigenvector transverse to $\{\rho=0\}$, and the other eigenvalues, those of the Riemannian Hessian, are negative. By the stable manifold theorem~\cite[Theorems~9.4 and~9.5]{Tes12}, the points whose backward orbit stays in a small neighbourhood $V$ of $(0,u^*)$ form a $C^1$ curve $\Gamma$ transverse to $\{\rho=0\}$, and these orbits converge to $(0,u^*)$ as $\tau\to-\infty$. So $\Gamma\cap\{\rho>0\}$ is an arc of a single orbit. Since $dt=\rho^{2-k}d\tau$, it gives a gradient-flow trajectory $\gamma$ in $W^c$ with $\gamma\to\theta^*$ and $\gamma_n/|\gamma_n|\to u^*$ as $t\to-\infty$. Conversely,
\begin{equation}\tag{$\star$}\label{eq:star}
\parbox{0.84\textwidth}{every gradient-flow trajectory in $S_{A'}$ with $0<\rho\le\rho_1$, $|u-u^*|\le r$ and $|a-\theta^*|\le r$ for all $t\le t'$ is a time translate of $\gamma$.}
\end{equation}
Indeed, it lies in $W^c$, and $\kappa\rho^{k-1}\le\dot\rho\le C\rho^{k-1}$ gives $\rho\to0$ and $\tau\to-\infty$ as $t\to-\infty$. So $(\rho,u)$ is a backward orbit of $X$ in $V$, and it lies in $\Gamma$.

\emph{Proof of \textup{(a)}.} Existence and uniqueness follow, since every escape orbit of $\theta^*$ along $(n,u^*)$ satisfies the assumptions of~\eqref{eq:star} for $t'\ll0$. As $h(w)-\theta^*=o(|w|)$, the invariance step applies to $\gamma$ at every $t_0\ll0$, with $\eta=0$. So $|\gamma_n|$ increases until it reaches $\rho_1$, at $t=0$ say, and $|\pi_A\gamma-\theta^*|\le C|\gamma_n|^k$ by~\eqref{eq:slaving} with $t_0\to-\infty$. The last claim is a special case of (b).

\emph{Shadowing: proof of \textup{(b)}.} The first two claims are contained in the invariance step. For the others, let $\theta_j$ be trajectories as in (b) with $\rho_0\to0$, shifted so that $t_1=0$. Then $t_0\to-\infty$, since $\dot\rho\le C\rho^{k-1}$. By compactness and continuous dependence on initial conditions, a subsequence converges, locally uniformly on $(-\infty,0]$, to a gradient-flow trajectory $\theta_\infty$. By the invariance step, $\theta_\infty$ satisfies the assumptions of~\eqref{eq:star} with $t'=0$, and its $\rho$ increases to $\rho_1$ at $t=0$. Hence $\theta_\infty=\gamma$, for every subsequence, and $\theta_j\to\gamma$ uniformly on every $[-T,0]$. Finally, by~\eqref{eq:slaving}, $\theta_j(t)$ is $O(\rho')$-close to $\theta^*$, a limit point of $\gamma$, as long as $\rho\le\rho'$, and $\rho\ge\rho'$ only for $t\ge-C\rho'^{2-k}$. Let $j\to\infty$ and then $\rho'\to0$.
\end{proof}

\subsection*{Proof of Theorem~\ref{thm:main}}

We prove by induction on $m\le M$:

\smallskip\noindent$(\mathcal C_m)$\quad there are times $\tau^\alpha_m$ with $\alpha^{k-2}\tau^\alpha_m\to s_m$, $|\pi_{A_m}\theta^\alpha(\tau^\alpha_m)-\theta^{(m)}|\le r_0$ (with $r_0$ from Lemma~\ref{lem:plateau} at $\theta^{(m)}$), and $\alpha^{-1}w^\alpha_n(\tau^\alpha_m)\to\Xi_n(s_m)$ for $n\notin A_m$.

\smallskip\noindent$(\mathcal C_0)$ holds with $\tau^\alpha_0=0$.

\emph{Plateau.} Assume $(\mathcal C_m)$. If $m<M$, apply Lemma~\ref{lem:plateau} on $[s_m,s_{m+1}-\delta]$; if $m=M$, apply it on $[s_M,S']$ with $S<S'<s_{M+1}$. This gives (ii) on this interval. Since the inactive weights are $O(\alpha)$, it also gives $|\theta^\alpha-\theta^{(m)}|\le e^{-\mu(t-\tau^\alpha_m)}r_0+O(\alpha)$, which tends to $0$ uniformly for $s\ge s_m+\delta$. This is (i). For $m=0$, (i) holds on the whole interval, since $\pi_\emptyset\theta=0$.

\emph{Jump.} Let $m<M$, $n_*=n_{m+1}$ and $u_*=u_{m+1}$, and take $\rho_1,r,\eps_1$ from Lemma~\ref{lem:escape}(b) at $(A_m,\theta^{(m)},n_*,u_*)$. By \ND2, $|\Xi_{n_*}(s)|\to\infty$ with direction converging to $u_*$ as $s\uparrow s_{m+1}$. By \ND1, the other dormant states stay bounded on $[s_m,s_{m+1}]$: $\sum_{n\notin A_{m+1}}|\Xi_n|^2\le R'^2$. Fix $\delta>0$ small, so that $R_\delta:=|\Xi_{n_*}(s_{m+1}-\delta)|\ge2R'/\eps_1$ and $\Xi_{n_*}(s_{m+1}-\delta)/R_\delta$ is $r/2$-close to $u_*$. Let $t_0$ be the time corresponding to $s_{m+1}-\delta$. For small $\alpha$, the assumptions of Lemma~\ref{lem:escape}(b) hold at $t_0$. Indeed, by Lemma~\ref{lem:plateau}, $w_{n_*}=\alpha(\Xi_{n_*}(s_{m+1}-\delta)+o(1))$, so $\rho_0=\alpha(R_\delta+o(1))$ and $|u-u_*|\le r$. Moreover, $\eta\le\alpha(R'+o(1))\le\eps_1\rho_0$ and $|a-\theta^{(m)}|=O(\alpha^k)\le\rho_0$. Let $t^\alpha_{m+1}:=t_1$ be the first time after $\tau^\alpha_m$ at which $|w_{n_*}|=\rho_1$; it does not depend on $\delta$. Lemma~\ref{lem:escape}(b) gives:
\begin{itemize}[nosep]
\item $\theta^\alpha(t_1)\to\gamma_{m+1}(0)$, where $\gamma_{m+1}$ is normalized as in the lemma, since $\rho_0\to0$;
\item the jump is fast: $\alpha^{k-2}(t_1-t_0)\le C\alpha^{k-2}\rho_0^{2-k}\to CR_\delta^{2-k}$;
\item $|z|\le2\eta=O(\alpha)$ on $[t_0,t_1]$. So $|\frac{d}{ds}(\alpha^{-1}w_n)|=\alpha^{1-k}|\nabla_{w_n}F|=O(1)$ by Lemma~\ref{lem:flat}(a), and the other dormant neurons move by $O(R_\delta^{2-k})$ in rescaled units, uniformly in $\delta$.
\end{itemize}
By continuous dependence on initial conditions, forward and backward in time, $\theta^\alpha(t_1+\cdot)\to\gamma_{m+1}$ uniformly on compact intervals. By \ND3, there is $T$ with $|\gamma_{m+1}(T)-\theta^{(m+1)}|\le r_0/2$. Set $\tau^\alpha_{m+1}=t_1+T$. On $[t_1,t_1+T]$, the dormant weights move by $O(\alpha^{k-1})$. Hence $\alpha^{k-2}\tau^\alpha_{m+1}\in[s_{m+1}-\delta,\,s_{m+1}-\delta+CR_\delta^{2-k}]+o(1)$. Moreover, for $n\notin A_{m+1}$,
\[
\big|\alpha^{-1}w_n^\alpha(\tau^\alpha_{m+1})-\Xi_n(s_{m+1})\big|\le CR_\delta^{2-k}+\operatorname{osc}_{[s_{m+1}-\delta,s_{m+1}]}\Xi_n+o(1).
\]
Since $\tau^\alpha_{m+1}$ does not depend on $\delta$, and $R_\delta\to\infty$ as $\delta\to0$, this proves $(\mathcal C_{m+1})$, for every large $T$.

\emph{Conclusion.} The plateau and jump steps cover $[0,\alpha^{2-k}S]$, and $\theta^\alpha$ stays bounded there. This gives the existence of $\theta^\alpha$ on this interval, (i), (ii), and the local uniform convergence in (iii). For the Hausdorff convergence at $s_{m+1}$, run the jump step with some $\delta'<\delta$. The curve is close to $\theta^{(m)}$ before $t_0$, by (i). On $[t_0,t_1]$ it is close to $\gamma_{m+1}((-\infty,0])$, by Lemma~\ref{lem:escape}(b), and on $[t_1,t_1+T]$ to $\gamma_{m+1}([0,T])$. Afterwards it stays within $|\gamma_{m+1}(T)-\theta^{(m+1)}|+o(1)$ of $\theta^{(m+1)}$, by Lemma~\ref{lem:plateau}, and $T$ is arbitrarily large. The converse follows from the local uniform convergence and from (i) at $s_{m+1}\pm\delta$. Finally, $F(\theta^{(m)})>F(\theta^{(m+1)})$ because $F$ decreases strictly along $\gamma_{m+1}$. For $m<M$, $\theta^{(m)}$ is a saddle by the consequences of Lemma~\ref{lem:flat}, since $U^{(m)}_{n_{m+1}}(u_{m+1})>0$.\qed

\begin{thebibliography}{9}\small

\bibitem{Car81} J.~Carr. \emph{Applications of Centre Manifold Theory}. Applied Mathematical Sciences~35, Springer, 1981.

\bibitem{KMC25} D.~Kunin, G.~L.~Marchetti, F.~Chen, D.~Karkada, J.~B.~Simon, M.~R.~DeWeese, S.~Ganguli, N.~Miolane. Alternating gradient flows: a theory of feature learning in two-layer neural networks. NeurIPS, 2025. arXiv:2506.06489.

\bibitem{Tes12} G.~Teschl. \emph{Ordinary Differential Equations and Dynamical Systems}. Graduate Studies in Mathematics~140, American Mathematical Society, 2012.

\end{thebibliography}

\end{document}
